\chapter{TODO: Title of Study One}\label{chapter:study-one}
\thispagestyle{myheadings}

	% ========================================================================
	% RESULTS CHAPTER
	%
	% One chapter per study, each roughly corresponding to one paper. The usual
	% shape is: a short introduction that narrows from the general introduction
	% to this specific question, any methods unique to this study, the results
	% organised so that each section makes one point, and a short discussion
	% that is picked up again in the final chapter.
	%
	% If this chapter is published or under review, note it here AND on the
	% copyright page (BU guide 1.3.2).
	% ========================================================================

	\section{Introduction}\label{section:study-one:introduction}

		TODO: Two to four paragraphs. Narrow from the broad framing of
		\cref{chapter:introduction} to the specific question of this chapter,
		and end with a one-sentence statement of what you found.

	\section{Additional methods}\label{section:study-one:methods}

		TODO: Only what is specific to this study. Everything shared belongs in
		\cref{chapter:methods}; cross-reference it rather than repeating it,
		for example "imaging was performed as described in
		\cref{section:methods:optics}".

	\section{Results}\label{section:study-one:results}

		\subsection{TODO: first finding as a declarative sentence}

			TODO: Section headings that state the finding ("Stimulation
			increased X by Y") read far better than headings that name the
			measurement ("Analysis of X").

			% ----------------------------------------------------------------
			% Figures. Keep each figure in its own file under figures/ and the
			% artwork under graphics/, then \input the wrapper here -- that is
			% how the rest of this template is organised.
			%
			% BU guide 1.1.8: number figures and tables uniquely, either
			% sequentially through the whole document or by chapter (3.1, 3.2).
			% Do NOT restart numbering in each chapter.
			%
			% BU guide 1.11: make sure everything is legible, avoid dark
			% backgrounds with dark text, and split a figure across pages
			% rather than shrinking it until it cannot be read.
			%
			% The optional argument to \caption is what appears in the List of
			% Figures -- keep it short.
			% ----------------------------------------------------------------
			% \begin{figure}[htbp]
			% 	\centering
			% 	\includegraphics[width=\textwidth]{your-figure}
			% 	\caption[Short caption for the List of Figures]{
			% 		Full caption. Describe what is plotted, what the error
			% 		bars represent, the n for animals and cells, and the
			% 		statistical test. A caption should be readable on its own.
			% 	}\label{figure:study-one:example}
			% \end{figure}

		\subsection{TODO: second finding}

			TODO.

	\section{Discussion}\label{section:study-one:discussion}

		TODO: Interpret this study's findings specifically. Save the synthesis
		across studies for \cref{chapter:discussion}.
